\documentclass[11pt,a4paper]{article}

\usepackage[T1]{fontenc}
\usepackage[utf8]{inputenc}
\usepackage[british]{babel}
\usepackage{amsmath,amssymb}
\usepackage{booktabs}
\usepackage{csquotes}
\usepackage{float}
\usepackage{geometry}
\usepackage{microtype}
\usepackage{tabularx}
\usepackage{xcolor}
\usepackage[hidelinks]{hyperref}
\usepackage[backend=biber,style=authoryear,maxcitenames=2,maxbibnames=6,
doi=true,isbn=false,url=false]{biblatex}

\geometry{margin=24mm}
\addbibresource{references.bib}
\setlength{\parindent}{0pt}
\setlength{\parskip}{0.55em}
\renewcommand{\arraystretch}{1.12}
\definecolor{accent}{HTML}{274C77}
\hypersetup{
  colorlinks=true,
  linkcolor=accent,
  citecolor=accent,
  urlcolor=accent,
  pdftitle={Duration of antibody protection project},
  pdfsubject={Nagoya-Oxford workshop 2026 participant handout}
}

\title{\textbf{Duration of antibody protection project}\\[0.35em]
\large Nagoya-Oxford workshop 2026}
\author{}
\date{August 2026}

\begin{document}
\maketitle
\vspace{-2.2em}

\section{Project question and aims}\label{sec:aims}

Consider a vaccine trial, or its long-term extension, in which an
antibody marker is measured repeatedly after vaccination. Antibody
levels often fall quickly at first and then more slowly. We want to
estimate when the marker will fall below a level associated with
protection, even if that point lies beyond the period observed so far.
This creates a practical decision at each planned trial review:

\begin{quote}
  \emph{Given the number and timing of the measurements collected so
    far, how well can we estimate the duration of antibody-associated
  protection, and how much would we gain from further follow-up?}
\end{quote}

The project has two core aims, followed by two extensions:
\begin{enumerate}
  \item under a known trajectory model, determine how well its
    parameters and the duration of protection can be identified as the
    length and frequency of follow-up vary;
  \item when the trajectory model is unknown, quantify the additional
    uncertainty from considering several plausible long-term patterns;
  \item examine how differences between participants affect these
    results and the balance between recruiting more participants and
    measuring each participant more often;
  \item turn the findings into practical guidance on how long, and how
    often, participants should be followed.
\end{enumerate}

Ending follow-up for this durability question is not the same as
ending all trial activity. Safety surveillance, clinical endpoints and
regulatory commitments may require the trial to continue for other
reasons. This project asks only when enough information has been
collected to answer the durability question.

The problem is most apparent when protection is expected to last for
years. The pivotal Shingrix trials followed participants for roughly
three to four years, and their long-term extension ultimately assessed
efficacy and immunogenicity up to 11 years after vaccination
\parencite{strezova2025}. Such studies show the value of long-term
follow-up. They also raise a practical question: could the antibody
data available at an earlier review already have supported a reliable
conclusion, or were later measurements needed to establish the
long-term pattern?

Protection after COVID-19 vaccination or infection usually wanes over
months rather than years. Even so, COVID-19 is a useful first case
study: neutralisation predicts protection, and published studies link
antibody titres with variants, boosting and vaccine effectiveness
\parencite{khoury2021,cromer2022,gilbert2022,feng2021}. We can
therefore develop and test the methods in a comparatively rich setting
before potentially applying them to pathogens for which protection may last for
years but the data, or the evidence for an immune correlate, are more
limited.

\section{Defining duration of protection}\label{sec:estimand}

Let $c_i(t)$ denote the underlying level of the antibody marker for
participant $i$ at time $t$. The marker could be a neutralisation
titre, a binding-antibody concentration or another assay measurement.
Let $c_{\mathrm{thr}}$ denote a protection-relevant threshold on the
same scale. We use the first time the marker crosses this threshold as
an indication of the participant's duration of protection:
\begin{equation}
  T_i^* = \inf\{t\geq 0 : c_i(t)\leq c_{\mathrm{thr}}\}.
  \label{eq:tstar}
\end{equation}
We note that this quantity depends on how protection and the
threshold are defined, and
that it has no finite value if the antibody level settles above the
threshold.

In the initial COVID-19 application, $c_i(t)$ and $c_{\mathrm{thr}}$
refer specifically to neutralising-antibody titres. Let
$c_{\mathrm{conv}}$ be the mean neutralisation level among
convalescent people, and divide titres by this reference, so that the
mean convalescent level is 1. One option is to set
$c_{\mathrm{thr}}$ to the relative titre that
\textcite{khoury2021} associated with 50\% protection. Alternatively,
we could use their full relationship between titre and protection to
convert the fitted antibody trajectory into a protection curve
$E_i(t)$. The latter approach makes clear that antibody persistence
and protection are related but are not the same thing. Antibody
markers can provide useful quantitative predictions of protection
even though they capture only part of the underlying immune response
\parencite{plotkin2012,earle2021}.

Differences between participants create a further distinction. The
distribution of their individual $T_i^*$ values is not the same as
the crossing time calculated from an average antibody curve. A useful
population measure might instead be the time at which the proportion
of people still above the threshold falls below a chosen value. A
later stage of the project will examine both the individual and
population questions.

\section{Proposed analysis programme}\label{sec:analysis}

The central idea is to simulate complete antibody histories for a set
of participants and recreate the information that would have been
available at a series of trial reviews. At each review, we would hide
all later measurements, fit one or more models to the data collected
by that point, and estimate $T^*$. Because the complete simulated
histories and generating models are known, we can first isolate how
much information the trial design provides and then quantify the
additional uncertainty that arises when the correct model is unknown.

The workshop should first agree on the clinically relevant quantity
and on what result would justify ending durability follow-up. This
requires an acceptable level of uncertainty, an acceptable risk of a
wrong decision, and a limit on extrapolation beyond the observed data.
We should also decide whether COVID-19 is primarily a worked example
or the main application.

The work could then proceed in three parts:
\begin{enumerate}
  \item \textbf{Known trajectory model: parameter identifiability.}
    Begin with a biphasic model used both to generate and fit the data.
    Vary the review date, visit schedule, sample size, assay noise and
    quantification limits. At each review, assess how well the model
    parameters and $T^*$ are identified. This isolates the information
    provided by the amount and pattern of observation.

  \item \textbf{Unknown trajectory model: model uncertainty.}
    Generate data from each plausible trajectory in turn, and later from
    one of the more detailed biological models. Fit the full set of
    candidate models without assuming which is correct. At each review,
    separate the uncertainty within each fitted model from the
    uncertainty arising from the choice of model, and assess how both
    affect estimation of $T^*$.

  \item \textbf{Extensions and practical interpretation.} Allow
    participants to differ in their peak antibody levels and rates of
    decline; compare individual and population outcomes; and compare
    the value of measuring more participants with that of measuring each
    participant more often. If useful, convert antibody trajectories
    into protection against symptomatic and severe disease, and apply
    the same workflow to a suitable open dataset.
\end{enumerate}

The main output would be a set of conditions under which estimates of
duration are reliable enough to inform follow-up decisions. These
conditions would distinguish limitations caused by insufficient data
under a known model from those caused by uncertainty about the model
itself.

\section{Models of antibody decline}\label{sec:models}

The analysis requires two kinds of model. Relatively simple models
will be fitted to the measurements available at each trial review.
More detailed biological models can then generate test data with
features that the fitted models may not capture. Keeping these roles
separate makes it possible to ask whether a simple analysis remains
useful when its assumptions are not exactly true.

\subsection*{Models to fit}

Each fitted model makes a different assumption about what happens
after follow-up ends. We suggest using the biphasic exponential as the
reference for the first analyses, while keeping the other models as
credible alternatives.

\subsubsection*{Single exponential}
\begin{equation}
  c(t)=c_0e^{-kt}.
\end{equation}
This is the simplest model. The log concentration falls at a constant
rate, and $T^*=k^{-1}\log(c_0/c_{\mathrm{thr}})$ when
$c_0>c_{\mathrm{thr}}$. If antibody levels really do follow this
pattern, relatively short follow-up may be enough to estimate the
curve. If they do not, the model wrongly extends the initial rate of
decline indefinitely. Describing COVID-19 neutralisation in terms of a
single half-life is nevertheless common
\parencite{khoury2021,cromer2022}.

\subsubsection*{Biphasic exponential (suggested default)}
Biphasic models describe an initial rapid fall followed by a slower,
more persistent decline, and have been used successfully to model
antibody responses after vaccination \parencite{white2015,hogan2023}.
We suggest this family as the initial reference because its parameters
have useful biological interpretations, the part that governs
long-term behaviour is explicit, and published COVID-19 estimates are
available for calibration. It is a practical basis for comparison,
not a claim that biphasic decay will suit every dataset. Following
\textcite{hogan2023}, we write the model as
\begin{equation}
  c(t)=c_0
  \frac{
    \exp\!\left[-(\log 2)\left(\frac{t}{h_1}+\frac{t_s}{h_2}\right)\right]
    +\exp\!\left[-(\log 2)\left(\frac{t}{h_2}+\frac{t_s}{h_1}\right)\right]
  }{
    \exp[-(\log 2)t_s/h_1]+\exp[-(\log 2)t_s/h_2]
  },
  \qquad h_1<h_2.
  \label{eq:biphasic}
\end{equation}
The fast component can represent the loss of short-lived
antibody-secreting cells, and the slow component a longer-lived cell
population. This parameterisation ensures that $c(0)=c_0$ and that the
two components contribute equally at the switching time $t_s$. The slow
half-life $h_2$ determines the long-term behaviour, but may be hard to
estimate until the fast component has subsided.

\subsubsection*{Power-law decay}
\begin{equation}
  c(t)=c_0(1+\alpha t)^{-\beta}.
\end{equation}
Here the rate of decline slows continuously. Over a short period the
curve can look exponential, yet it predicts that antibodies will
persist much longer. Long-term antibody measurements give biological
grounds for considering this pattern rather than assuming a constant
half-life \parencite{amanna2007}.

\subsubsection*{Plateau or set-point}
\begin{equation}
  c(t)=A_\infty +(c_0-A_\infty)e^{-kt}.
\end{equation}
In this model, the antibody level settles at a non-zero value. It
crosses the threshold only if $A_\infty<c_{\mathrm{thr}}$. With short
follow-up, a true plateau may be indistinguishable from a very slow
decline. Long-term
measurements showing persistent antibody levels motivate this option
\parencite{amanna2007}. Because $T^*$ may not have a finite value under
this model, the workshop should decide whether and how to include it in
model comparisons.

\subsection*{More complex models for generating test data}

\subsubsection*{Mechanistic antibody-production model}
A more detailed model can represent short- and long-lived
antibody-secreting cells, as well as the clearance of the antibodies
themselves:
\begin{align}
  \frac{dP_s}{dt}&=-\delta_sP_s, &
  \frac{dP_l}{dt}&=-\delta_lP_l, &
  \frac{dc}{dt}&=\rho_sP_s+\rho_lP_l-\delta_cc.
\end{align}
When antibodies clear more quickly than the cells, this model behaves
like biphasic decay. Adding immune memory or cell replenishment can
produce still slower long-term decline. Related models have been
developed for plasma-cell dynamics and vaccine responses
\parencite{andraud2012,pasin2019}. For this project, such a model may
be most useful for generating realistic test data, rather than as the
default model fitted to every dataset.

\subsubsection*{COVID-19 semi-mechanistic model}
A second option for generating test data is the semi-mechanistic model
developed in a series of COVID-19 studies
\parencite{nishiyama2023,nakamura2024,park2025,hart2025}. It links a
waning mRNA stimulus to delayed, non-linear antibody production and
antibody loss, and can account for boosting, infection history and
differences between participants. Once the vaccine stimulus has
disappeared, however, the model approaches a single exponential
decline.

\section{Measurement error and starting values for the biphasic
model}\label{sec:parameterisation}

Assay results are noisy measurements of an underlying positive
trajectory. A simple starting point is to assume multiplicative
log-normal error:
\begin{equation}
  \log y_{ij}=\log c_i(t_{ij})+\epsilon_{ij},
  \qquad \epsilon_{ij}\sim N(0,\sigma_{\log}^2).
  \label{eq:observation}
\end{equation}
Results below or above the assay's quantification limits can be treated
as censored, rather than assigned exactly to the relevant limit. This
observation model can be combined with any of the trajectories in
Section~\ref{sec:models}.

Table~\ref{tab:starting-values} gives possible starting parameters for
a biphasic exponential model of BNT162b2 after dose 3 against the
Delta variant, based on the parameterisation of \textcite{hogan2023}.
All antibody levels are expressed relative to the mean convalescent
neutralisation level defined above.

\begin{table}[H]
  \centering
  \footnotesize
  \caption{Tentative starting values for a BNT162b2 dose-3 biphasic
  COVID-19 neutralising-antibody baseline.}
  \label{tab:starting-values}
  \begin{tabularx}{\textwidth}{@{}
      >{\raggedright\arraybackslash}p{0.22\textwidth}
      >{\raggedright\arraybackslash}p{0.16\textwidth}
      >{\raggedright\arraybackslash}p{0.23\textwidth}
    >{\raggedright\arraybackslash}X@{}}
    \toprule
    Quantity & Starting value & Source / basis & Interpretation /
    possible variation \\
    \midrule
    Early level $c_0$ & 0.92 &
    \textcite{hogan2023} & BNT162b2 dose 3 against Delta; 95\% CrI 0.71--1.19 \\
    Log-error SD $\sigma_{\log}$ & 0.20 & Workshop noise choice &
    Explore lower and higher assay variability \\
    Fast half-life $h_1$ & 35 days & \textcite{hogan2023} & Early
    contraction; tentative range 25--60 days \\
    Slow half-life $h_2$ & 581 days & \textcite{hogan2023} &
    Most important long-term assumption; consider 250--900 days \\
    Switching time $t_s$ & 75 days & \textcite{hogan2023} & Equal
    contributions from the two components at this time \\
    Protection anchor $c_{\mathrm{thr}}$ & 0.202 / 0.030 &
    \textcite{khoury2021} & Relative titres associated with 50\%
    symptomatic / severe protection \\
    \bottomrule
  \end{tabularx}
\end{table}

\section{Differences between participants}\label{sec:heterogeneity}

In the biphasic reference model, participants can differ in their peak
antibody level and slow half-life. One way to represent these
differences is with correlated log-normal random effects directly on
those two quantities:
\begin{align}
  c_{0i} &= c_0e^{\eta_{0i}},
  & h_{2i} &= h_2e^{\eta_{hi}}, \\
  \boldsymbol{\eta}_i
  &=
  \begin{pmatrix}\eta_{0i}\\ \eta_{hi}
  \end{pmatrix}
  \sim N_2(\boldsymbol{0},\boldsymbol{\Sigma}), \qquad
  & \boldsymbol{\Sigma}
  &=
  \begin{pmatrix}
    \omega_0^2 & \rho_{0h}\omega_0\omega_h \\
    \rho_{0h}\omega_0\omega_h & \omega_h^2
  \end{pmatrix}.
  \label{eq:heterogeneity}
\end{align}
This structure lets us estimate outcomes for both individuals and the
population. It also lets us ask whether recruiting more participants
can make up for taking only a few measurements from each person.
Similar structures have been used in population models of vaccine
responses \parencite{pasin2019}.

Table~\ref{tab:heterogeneity} gives tentative starting values for this
between-participant variation, using the same BNT162b2 dose-3 setting
as Table~\ref{tab:starting-values}.

\begin{table}[H]
  \centering
  \footnotesize
  \caption{Tentative starting values for between-participant
  heterogeneity in the biphasic baseline.}
  \label{tab:heterogeneity}
  \begin{tabularx}{\textwidth}{@{}
      >{\raggedright\arraybackslash}p{0.22\textwidth}
      >{\raggedright\arraybackslash}p{0.16\textwidth}
      >{\raggedright\arraybackslash}p{0.23\textwidth}
    >{\raggedright\arraybackslash}X@{}}
    \toprule
    Quantity & Starting value & Source / basis & Interpretation /
    possible variation \\
    \midrule
    Peak / half-life log-SD $\omega_0/\omega_h$ & 0.30 / 0.25 & Workshop
    starting value & Moderate differences between participants \\
    Peak--half-life correlation $\rho_{0h}$ & $0.30$ & Workshop
    correlation choice & Equivalent to a $-0.30$ peak--rate
    correlation; include zero as a comparison \\
    \bottomrule
  \end{tabularx}
\end{table}

The uncited values in Tables~\ref{tab:starting-values} and
\ref{tab:heterogeneity} are convenient starting points for simulation,
not estimates from the literature. The workshop may instead choose a
different vaccine, dose, assay scale or long-term rate of decline.

\section{Evidence base and key studies}\label{sec:literature-guide}

The papers below are grouped by the part of the project they inform:
interpreting antibody markers, calibrating the COVID-19 example,
representing differences between participants, and understanding
long-term antibody dynamics and follow-up.

\subsection*{From antibody markers to protection}

\textbf{Plotkin and Gilbert (2012).} This framework distinguishes a
statistical correlate, a mechanistic correlate and a surrogate
endpoint. It helps us avoid treating a protection-relevant antibody
marker as the whole immune mechanism
\parencite{plotkin2012}.

\textbf{Earle et al. (2021).} This comparison across vaccines gives
broader evidence that antibody levels can act as a correlate of
protection for COVID-19 vaccines. It supports our use of an antibody
marker without implying that antibodies alone cause protection
\parencite{earle2021}.

\textbf{Gilbert et al. (2022).} This analysis of the mRNA-1273
efficacy trial shows that neutralisation is a graded correlate of
risk. It therefore supports modelling the marker continuously rather
than using only a binary laboratory cut-off \parencite{gilbert2022}.

\textbf{Feng et al. (2021).} This analysis of ChAdOx1 nCoV-19 also
finds a continuous relationship between immune markers and symptomatic
or asymptomatic infection. It cautions against treating one threshold
as universal \parencite{feng2021}.

\textbf{Khoury et al. (2021).} This study provides the main
quantitative link between neutralisation titre and protection. Its
logistic relationship and estimates of the titres associated with
50\% protection can be used to convert a fitted antibody trajectory
into a protection curve for a specific clinical endpoint
\parencite{khoury2021}.

\textbf{Cromer et al. (2022).} This meta-analysis extends the link
between neutralisation titre and protection across variants and
boosting. It shows that the initial titre and antigenic change can
alter the duration of protection even when antibodies wane at the same
rate
\parencite{cromer2022}.

\subsection*{COVID-19 antibody kinetics and calibration}

\textbf{Hogan et al. (2023).} This modelling study estimates initial
neutralising-antibody levels for specific vaccines and doses, together
with fast and slow half-lives and a switching time. These estimates
provide the tentative COVID-19 calibration and show how antibody
dynamics can be linked to longer-term effectiveness against variants
\parencite{hogan2023}.

\textbf{Nishiyama et al. (2023).} This study extends an antibody
model based on mRNA stimulation to include a booster dose and previous
infection. It shows how repeated exposures can alter the duration of
the response and provides a compact semi-mechanistic model for later
simulations \parencite{nishiyama2023}.

\subsection*{Variation between participants}

\textbf{Nakamura et al. (2024).} This study combines frequent
measurements from a small group with sparse measurements from more than
2,400 participants to reconstruct individual trajectories after
primary vaccination. Its fitted parameter distributions are useful
for calibrating variation between participants, although the clinical
data are access-restricted
\parencite{nakamura2024}.

\textbf{Park et al. (2025).} This study groups antibody trajectories
after primary and booster vaccination into durable, vulnerable and
rapidly declining patterns. The link between rapid decline and earlier
subsequent infection shows that differences between participants have
clinical importance \parencite{park2025}.

\textbf{Hart et al. (2025).} This study links the same individual
antibody model to susceptibility and population outbreak risk. Its
open code and derived parameter estimates can be used directly for
simulation and calibration \parencite{hart2025}.

\textbf{Pasin et al. (2019).} This Ebola prime--boost study shows how
biological antibody dynamics and differences between participants can
be estimated together. It is therefore a useful guide for the stage
of the project that introduces participant variation
\parencite{pasin2019}.

\subsection*{Long-term dynamics and follow-up}

\textbf{Amanna et al. (2007).} Long-term measurements across antigens
show that antibody persistence varies greatly and need not follow one
constant half-life. This finding motivates the power-law and plateau
models, and the project's focus on long-term behaviour
\parencite{amanna2007}.

\textbf{Andraud et al. (2012).} These linked plasma-cell and antibody
models provide a clear biological basis for generating trajectories
over several time scales. They can be used to test when a simpler
decay curve is a good enough approximation \parencite{andraud2012}.

\textbf{White et al. (2015).} This analysis of the RTS,S malaria
vaccine is the closest published example of the proposed approach: it
uses a biphasic model of antibody decline and links that decline to the
duration of vaccine efficacy \parencite{white2015}.

\textbf{Strezova et al. (2025).} The final Shingrix analysis shows the
practical scale of this design problem. The extension of the pivotal
phase 3 trials followed more than 7,000 participants and assessed
protection for up to 11 years. It provides a concrete example in which
the value of more durability data must be weighed against the burden
of prolonged follow-up
\parencite{strezova2025}.

\printbibliography[title={Selected literature}]

\end{document}
